\section{Problema de Investigación}
 
\subsection{Planteamiento del Problema}

La rápida adopción de la Inteligencia Artificial Generativa y la arquitecturas RAG ha transformado el desarrollo de software actual mediante la integración dinámica de modelos de lenguaje y bases vectoriales. Además de estar impulsadas por una intensa presión de mercado que ha relegado los estándares de calidad y ha normalizado malas practicas como código desordenado y prompts mezclados con la lógica de negocio \cite{shao_are_2024}.

De continuar sin intervención derivará en una deuda técnica creciente que solo revelará fallos críticos en producción y disparará los costos de mantenimiento. Problemática que este proyecto propone prever caracterizando dichos antipatrones y sentando bases metodológicas para el futuro desarrollo de una herramienta proactiva de análisis estático basada en arboles de sintaxis abstracta.
 
\subsection{Formulación del Problema}
 
¿De qué manera la presencia de antipatrones arquitectónicos en aplicaciones basadas en IA Generativa y RAG afecta su mantenibilidad del software?
 
\subsection{Sistematización del Problema}
 
\begin{itemize}
    \item ¿Cuál es la taxonomía conceptual de antipatrones arquitectónicos específicos para sistemas GenAI/RAG de la literatura y la industria?
    \item ¿Cómo se pueden detectar automáticamente estos antipatrones en proyectos de GitHub mediante una herramienta de análisis estático basada en Árboles de Sintaxis Abstracta?
    \item ¿Cuál es la frecuencia de aparición de estos antipatrones en los repositorios analizados y qué correlación estadística existe con las métricas tradicionales de acoplamiento y complejidad sintáctica?
\end{itemize}